\section{Field theory $\boldsymbol{D}\mathbf{+1}$ dimensions}\label{sec:fieldtheory}

The Feynman rules for the correlation functions of the time-dependent
gauge field obtained in the previous section
are those of a local field theory in $D+1$ dimensions \cite{ZinnZwanziger}.
It is possible to show this through somewhat formal
functional-integral mani\-pulations, but one can also adopt
a purely algebraic point of view, where the action in $D+1$ dimensions
merely serves as a generating function for the Feynman rules.

\subsection{Action}

Except for the boundary condition \eqref{eq:2.3},
the field $B_{\mu}(t,x)$ will now be considered to be an
independent field.
One also needs to introduce a Lagrange-multiplier
field $L_{\mu}(t,x)=L_{\mu}^a(t,x)T^a$
with purely imaginary components.
The action of the theory in $D+1$ dimensions is then given by
\cite{ZinnZwanziger}
\begin{align}
   \Stot=S+\Sgf+\Sgh+\Sfl,
   \label{eq:4.1}\\[2.5pt]
   \Sfl=-2\int_0^{\infty}\rmd t\int\rmd^Dx\,
   \tr\bigl\{L_{\mu}(t,x)
   \bigl(\partial_tB_{\mu}-
         D_{\nu}G_{\nu\mu}-\alpha_0D_{\mu}\partial_{\nu}B_{\nu}
   \bigr)(t,x)\bigr\}.
   \label{eq:4.2}
\end{align}
In this framework, the flow equation \eqref{eq:2.4} coincides with
the field equation obtained by varying the action with respect to the
Lagrange-multiplier field. Note that the latter is not required to
satisfy any particular boundary conditions.

\subsection{Propagators}\label{ssec:propagators4}

The quadratic part of the action is the sum of
\begin{equation}
   \int_0^{\infty}\rmd t\int\rmd^Dx\, L^a_{\mu}(t,x)
   \bigl(
   \partial_tB^a_{\mu}-\partial_{\nu}\partial_{\nu}B^a_{\mu}
   -(\alpha_0-1)\partial_{\mu}\partial_{\nu}B^a_{\nu}
   \bigr)(t,x)
   \label{eq:4.3}
\end{equation}
and the quadratic part of the action in $D$ dimensions.
The bulk fields $B_{\mu}$ and $L_{\mu}$
thus couple to each other only, but there is also
an implicit coupling to the fundamental gauge field
through the boundary condition \eqref{eq:2.3}.
This complication can easily be overcome by substituting
\begin{equation}
   B_{\mu}(t,x)=\int\rmd^Dy\,K_t(x-y)_{\mu\nu}A_{\nu}(y)+b_{\mu}(t,x).
   \label{eq:4.4}
\end{equation}
The field $b_{\mu}$ then satisfies homogenous boundary conditions,
while the first term in eq.~\eqref{eq:4.4} solves the linearized flow equation
and thus drops out in the action \eqref{eq:4.3}.

It follows from these remarks that the propagators of
the fundamental gauge field and the ghost fields
coincide with the expressions given in the previous section.
All other two-point functions vanish except for
\begin{equation}
   \left.\bigl\langle
   b^a_{\mu}(t,x)L^b_{\nu}(s,y)
   \bigr\rangle\right|_{\rm leading\;order}
   =\delta^{ab}H(t,x;s,y)_{\mu\nu},
   \label{eq:4.5}
\end{equation}
which is determined by the field equation
\begin{equation}
   \bigl\{\delta_{\mu\rho}\partial_t-
   \delta_{\mu\rho}\partial_{\sigma}\partial_{\sigma}
   -(\alpha_0-1)\partial_{\mu}\partial_{\rho}\bigr\}
   H(t,x;s,y)_{\rho\nu}=\delta_{\mu\nu}\delta(t-s)\delta(x-y)
   \label{eq:4.6}
\end{equation}
and the boundary condition $\left.H_{\mu\nu}(t,x;s,y)\right|_{t=0,s>0}=0$.
The unique solution of these equations is
\begin{equation}
   H(t,x;s,y)_{\mu\nu}=\theta(t-s)K_{t-s}(x-y)_{\mu\nu}
   \label{eq:4.7}
\end{equation}
and the $bL$ propagator is thus seen to coincide with the flow
propagator.

The two-point functions involving the $B$ field
may finally be calculated by combining eq.~\eqref{eq:4.4} with the results
obtained so far. Since the $AL$ and the $bb$ propagators vanish, one
quickly finds that the $BL$ propagator is equal to the flow propagator
and that the $BB$ and the $AB$ propagators are
equal to the expressions given in the previous section.

\subsection{Vertices}

The theory described by the action \eqref{eq:4.1} has the vertices of the
theory in $D$ dimensions plus those deriving from the bulk action
$\Sfl$. Recalling eqs.~\eqref{eq:2.6} and \eqref{eq:2.13}, it is straightforward to show that
the $LB^2$ and $LB^3$ vertices generated by
the interaction part of the latter,
\begin{equation}
   \left.\Sfl\right|_{\rm interaction}=
   2\int_0^{\infty}\rmd t\int\rmd^Dx\,
   \tr\bigl\{L_{\mu}(t,x)R_{\mu}(t,x)\bigr\},
   \label{eq:4.8}
\end{equation}
coincide with the vertices $\vbulk{2,0}$ and $\vbulk{3,0}$.

All vertices and propagators of the
Feynman rules of the previous sections are thus recovered.
A somewhat unusual feature of these rules is the presence
of a field (the Lagrange multiplier) that propagates
only through its mixing with the other fields.
The algebraic structure of the Feynman rules
is however entirely standard and does not involve any
prescriptions apart from the ones deriving from the expansion of
the action in powers of the fields.

\begin{figure}[t]
\centerline{\includegraphics[width=6.5cm]{images/diagram8}}
\caption{Examples of diagrams with flow-line loops. All such diagrams
vanish as a consequence of the retarded nature of
the flow propagator and the rules of dimensional regularization
\cite{ZinnZwanziger}.
}
\label{fig:loops}
\end{figure}

\subsection{Flow lines and flow-line loops}\label{ssec:flowlines}

Flow lines represent the $BL$ propagator. They can start
and end at the flow vertices but are never attached
to an ordinary vertex. Outward- and inward-directed
external flow lines are external $B$ and $L$ lines,
respectively.

Diagrams with closed flow lines comply with these rules,
but should be absent if
a complete matching with the Feynman rules of
the previous section is to be achieved
(cf.~point (d) at the end of sect.~\ref{sec:perturbation}).
Such diagrams are in fact equal to zero \cite{ZinnZwanziger}.
The diagrams 1 and 2 in fig.~\ref{fig:loops}, for example, vanish because
dimensional regularization sets the momentum integral to zero.
If there are two or more vertices in the
loop, as in diagram 3, the time integrations
vanish, because the flow propagator is retarded and thus forces
the flow times at the vertices to be squeezed to a range of
measure zero (singularities in the time
coordinates are excluded in view of
the regularization of the momentum integral).

If a lattice regularization is used, the vanishing of diagram 2 may
not be guaranteed, but one can always include a ghost
field in the action that cancels the flow-line loops algebraically,
i.e.~at the level of the Feynman integrands \cite{Zinn}.
With dimensional regularization,
however, this device is not needed and is therefore omitted here.
